\documentclass[11pt,a4paper]{article}
\usepackage[margin=2.2cm]{geometry}
\usepackage{fontspec}
\setmainfont{Latin Modern Roman}
\setmonofont{Latin Modern Mono}
\usepackage[provide=*]{babel}
\babelprovide[import,main]{romanian}
\usepackage{xcolor}
\usepackage{booktabs,tabularx,array}
\usepackage{graphicx}
\usepackage{enumitem}
\usepackage{listings}
\usepackage[most]{tcolorbox}
\usepackage[hidelinks]{hyperref}
\usepackage{titlesec}
\usepackage{fancyhdr}
\usepackage{needspace}

\definecolor{accent}{HTML}{1F5F8B}
\definecolor{trap}{HTML}{B3261E}

\titleformat{\section}{\Large\bfseries\color{accent}}{}{0pt}{}[\vspace{0.2em}\color{accent}\titlerule]
\titlespacing*{\section}{0pt}{1.6em}{0.8em}
\setlist{itemsep=0.25em,topsep=0.3em}
\setlength{\parindent}{0pt}
\setlength{\parskip}{0.4em}

\pagestyle{fancy}
\fancyhf{}
\renewcommand{\headrulewidth}{0.3pt}
\fancyhead[L]{Administrarea rețelelor}
\fancyhead[R]{Curs 7 -- DNS}
\fancyfoot[C]{\thepage}

\lstset{basicstyle=\ttfamily\small,backgroundcolor=\color{black!4},frame=single,
  rulecolor=\color{black!20},xleftmargin=0.5em,xrightmargin=0.5em,
  columns=fullflexible,keepspaces=true,showstringspaces=false,breaklines=true}

\newtcolorbox{capcana}[1][Capcană de examen]{enhanced,colback=trap!5,colframe=trap,
  fonttitle=\bfseries,title={\,#1},boxrule=0.6pt,arc=2pt,left=6pt,right=6pt}

% imagine: se afișează doar dacă fișierul există
\newcommand{\img}[3][0.85]{\IfFileExists{#2}{\begin{center}
  \includegraphics[width=#1\linewidth]{#2}\\[-0.2em]{\small\itshape #3}\end{center}}{}}

\newcommand{\cod}[1]{\texttt{#1}}

\begin{document}

\begin{center}
  {\Huge\bfseries\color{accent} DNS -- Domain Name System}\\[0.4em]
  {\large Administrarea rețelelor \textbullet{} Curs 7 \textbullet{} Secțiunea VI}
\end{center}

\section{Ce este DNS}
\begin{itemize}
  \item Translatează un nume într-o adresă IP. Poate face și procesul invers, numit \textbf{Reverse DNS}.
  \item Este un sistem distribuit de servere și protocoale care gestionează baze de date cu nume de calculatoare, servere sau servicii din Internet.
  \item Fiecare nume poate primi o \textbf{adresă IP publică (DNS public)} sau o \textbf{adresă IP privată (server DNS privat într-o rețea privată)}.
  \item \textbf{Rolul principal}: să rezolve un nume de domeniu într-o adresă IP.
  \item Scop: utilizatorul nu mai trebuie să memoreze adrese IP.
  \item Clientul trimite un nume și primește \textbf{IP-ul asociat}, pe care apoi îl folosește pentru a accesa serviciul.
  \item Aplicații server: \textbf{bind, CoreDNS, UnboundDNS, PowerDNS, DNS Server (Windows Server)}.
  \item Folosește \textbf{portul 53}, de obicei pe \textbf{UDP}, dar funcționează și pe \textbf{TCP}.
  \item Structură ierarhică, arborescentă, pe 3 niveluri:
  \begin{itemize}
    \item \textbf{ROOT}: serverele de bază.
    \item \textbf{TLDs}: rezolvă ultimul nume de după ultimul punct.
    \item \textbf{Authoritative}: rezolvă domeniile obișnuite.
  \end{itemize}
\end{itemize}

\img[0.8]{Images/DNS_AR.png}{Ierarhia serverelor DNS (curs AR, slide 74)}
\img[0.62]{Images/DNS_INT.jpg}{Schema procesului DNS pe Internet}

\section{ROOT DNS Servers}
\begin{itemize}
  \item \textbf{Rol}: Root nu rezolvă domeniul complet, ci doar îți spune care server \textbf{TLD} se ocupă de ultima etichetă (de exemplu \cod{.ro}).
  \item Există \textbf{13 servere root} (a--m), fiecare cu \textbf{IP fix și nume de domeniu fix} (\cod{a.root-servers.net} \ldots{} \cod{m.root-servers.net}).
  \item Fiecare dintre cele 13 este implementat \textbf{redundant în mai multe zone geografice} (harta este pe root-servers.org).
  \item Fiecare device care se poate conecta la Internet are \textbf{o listă inițială cu serverele root}, inclusă în sistemul de operare.
  \item Serverele root sunt \textbf{nerecursive}.
  \item Operatori (tabelul de pe slide): a și j = Verisign; b = USC-ISI; c = Cogent; d = Univ.\ Maryland; e = NASA; f = ISC; g = US DoD; h = US Army; i = Netnod; k = RIPE NCC; l = ICANN; m = WIDE Project.
\end{itemize}

\img[0.72]{Images/map_root_servers.jpg}{Harta serverelor root}
\img[0.78]{Images/Tabel_Root_SV.png}{Tabelul serverelor root (slide 75)}

\section{TLD Servers}
\begin{itemize}
  \item Sunt pe \textbf{al doilea nivel după ROOT}. Rezolvă ultimul domeniu de după ultimul punct din URL.
  \item IP-urile lor sunt introduse în baza de date a serverelor ROOT.
  \item ROOT rezolvă doar TLD-ul și \textbf{deleagă} restul către ele.
  \item Categorii:
  \begin{itemize}
    \item \textbf{ccTLD} (country): cod de țară, de obicei 2 litere (\cod{.ro}, \cod{.de}, \cod{.hu}, \cod{.ch}).
    \item \textbf{Infrastructure}: \cod{.arpa}, primul domeniu creat, gestionat de IANA.
    \item \textbf{gTLD} (generic), gestionate de IANA: \cod{.com}, \cod{.org}, \cod{.net}, \cod{.edu}, \cod{.gov}, \cod{.mil}, \cod{.top}, \cod{.xyz}, \cod{.win}\ldots
    \item \textbf{grTLD} (generic restricted): \cod{.name}, \cod{.biz}.
    \item \textbf{sTLD} (sponsored): \cod{.aero}, \cod{.asia}, \cod{.coop}, \cod{.jobs}.
    \item \textbf{tTLD} (test): \cod{.test}.
    \item \textbf{Reserved}: pentru rețele interne, teste, aplicații private. \textbf{Nu sunt rezolvate de ROOT}: \cod{.example}, \cod{.invalid}, \cod{.local}, \cod{.localhost}, \cod{.onion}, \cod{.internal}, \cod{.alt}.
  \end{itemize}
  \item TLD-ul nu rezolvă tot URL-ul. \textbf{Deleagă} către serverul de nume (NS) al primului domeniu de după TLD, citind de la dreapta la stânga.
  \item NS-urile către care deleagă TLD-ul se numesc \textbf{servere autorizate (authoritative)}.
\end{itemize}

\Needspace{14\baselineskip}
\section{DNS Records (RR)}
Pentru fiecare domeniu sau subdomeniu se creează o \textbf{zonă} (similar cu o bază de date text) cu Resource Records.

\medskip
\renewcommand{\arraystretch}{1.35}
\begin{tabularx}{\linewidth}{@{}l l X@{}}
\toprule
\textbf{Record} & \textbf{Denumire} & \textbf{Ce face}\\
\midrule
\textbf{NS}    & Name Server & IP-ul unui server DNS autorizat pentru domeniu/subdomeniu\\
\textbf{A}     & Address Record & Adresă IPv4 pentru un host obișnuit\\
\textbf{AAAA}  & Address Record & La fel ca A, dar IPv6\\
\textbf{CNAME} & Canonical name record & Alias (alt nume) pentru un record deja definit, de obicei de tip A. Practic îl redenumim.\\
\textbf{MX}    & Mail exchange record & Listă de servere de mail pentru domeniu/subdomeniu. Serverul de mail este un host cu MTA (Mail Transfer Agent) pe \textbf{portul 25 (smtp)}\\
\textbf{SRV}   & Service locator & Locația (IP) unui serviciu\\
\textbf{TXT}   & Text record & Informații text generale, pentru aplicații utilizator\\
\textbf{SPF}   & Sender Policy Framework & Ca TXT, tot text, folosit în aplicații de mail\\
\textbf{SOA}   & „State of Authority'' (cum scrie slide-ul) & Probabil cel mai important record din zonă. Spune că serverul curent este autoritatea pentru zonă\\
\bottomrule
\end{tabularx}

\bigskip
\textbf{Sintaxa SRV:}
\begin{lstlisting}
_service._proto.name. ttl IN SRV priority weight port target
\end{lstlisting}
\textbf{Exemplu (slide):}
\begin{lstlisting}
_sip._tcp.example.com. 86400 IN SRV 20 0 5060 backupbox.example.com
\end{lstlisting}

\Needspace{18\baselineskip}
\section{SOA Record}
Record \textbf{obligatoriu} în orice zonă DNS.

\img[0.6]{Images/SOA_RECORD_CURS_AR.png}{Structura SOA și exemplul pentru cs.upt.ro (slide 78)}

\subsection*{Câmpuri}
\begin{itemize}
  \item \textbf{TTL}: perioada în secunde cu care serverul DNS superior (cel care a delegat) își actualizează cache-ul pentru zonă, în mod obișnuit, exceptând erorile.
  \item \textbf{Owner name}: zona. De obicei se scrie \cod{@}, care reprezintă direct numele zonei.
  \item \textbf{Class}: tipul zonei, de obicei \cod{IN} (Internet).
  \item \textbf{MNAME}: serverul NS \textbf{primar} al zonei, definit ulterior în zonă.
  \item \textbf{RNAME}: mailul administratorului, scris cu altă sintaxă:
  \begin{itemize}
    \item \textbf{primul punct} de la stânga la dreapta devine \cod{@};
    \item dacă vrei punct în partea de dinainte de \cod{@}, îl scrii cu escape \cod{\textbackslash.};
    \item \cod{admin@cs.upt.ro} $\rightarrow$ \cod{admin.cs.upt.ro};
    \item \cod{ion.popescu@cs.upt.ro} $\rightarrow$ \cod{ion\textbackslash.popescu.cs.upt.ro}.
  \end{itemize}
\end{itemize}

\textbf{Exemplu pentru \cod{cs.upt.ro}:}
\begin{lstlisting}
$TTL 86400
@   IN SOA   ns.cs.upt.ro. admin.cs.upt.ro. (
         2020080302 ;Serial
         7200       ;Refresh
         3600       ;Retry
         1209600    ;Expire
         3600       ;Negative response caching TTL
)
\end{lstlisting}

\begin{itemize}
  \item \textbf{Serial Number}:
  \begin{itemize}
    \item se schimbă la fiecare modificare a zonei;
    \item format \textbf{AAAALLZZNN}: AAAALLZZ = data modificării (an, lună, zi), NN = număr incrementat la fiecare modificare din aceeași zi;
    \item identifică modificările pe serverul primar (master). Când un \textbf{secondary} vede că serialul s-a schimbat, pornește actualizarea zonei prin \textbf{zone transfer}.
  \end{itemize}
  \item \textbf{Refresh}: la cât timp (secunde) citesc serverele secundare zona de la primar.
  \item \textbf{Retry}: după cât timp (secunde) reîncearcă secundarele dacă primarul nu răspunde.
  \item \textbf{Expire}: timeout. După această perioadă de inactivitate a primarului, secundarele \textbf{nu mai răspund la interogări}.
  \item \textbf{Minimum TTL}: similar cu TTL-ul inițial. TTL-ul final $=$ \textbf{minim(TTL, MINIMUM TTL)}.
\end{itemize}

\section{SOA cu configurație minimă}
\begin{lstlisting}
      IN NS    ns.cs.upt.ro.
      IN NS    ns.upt.ro.
      IN NS    ns1.upt.ro
      IN A     193.226.12.8
ns    IN A     193.226.12.13
ns2   IN A     193.226.12.3
web   IN A     193.226.12.8
www   IN CNAME web
webmail IN CNAME www
\end{lstlisting}
\begin{itemize}
  \item Zona are 3 NS-uri, un A pentru domeniul însuși (fără nume în față) și A-uri pentru \cod{ns}, \cod{ns2} și \cod{web}.
  \item \cod{www} este alias pentru \cod{web}, iar \cod{webmail} este alias pentru \cod{www}.
\end{itemize}

\img[0.8]{Images/Arhitectura_DNS_infografic.jpg}{Infografic: arhitectura și funcționarea sistemului DNS}

\begin{capcana}
\begin{itemize}
  \item Rootul \textbf{nu rezolvă domenii complete}: se oprește la TLD și este \textbf{nerecursiv}.
  \item Cel mai din dreapta din URL este \textbf{TLD-ul}. Rootul nu apare în URL.
  \item \textbf{13} servere root (a--m), dar implementate redundant geografic.
  \item Portul 53 este \textbf{de obicei UDP}, dar merge și pe TCP. Portul 25 este pentru MX/SMTP.
  \item Domeniile \textbf{Reserved} (\cod{.local}, \cod{.onion}, \cod{.example}\ldots) \textbf{nu sunt rezolvate de ROOT}; nu le confunda cu \cod{.test} (tTLD).
  \item \cod{.arpa} este \textbf{infrastructure}, nu gTLD. \textbf{grTLD} (\cod{.name}, \cod{.biz}) $\neq$ \textbf{sTLD} (\cod{.aero}, \cod{.asia}, \cod{.coop}, \cod{.jobs}).
  \item \textbf{CNAME} este alias pentru un record existent, de obicei A.
  \item \textbf{SOA este obligatoriu} în orice zonă.
  \item \textbf{RNAME}: primul punct $=$ \cod{@}; punctul din numele dinainte de \cod{@} se scrie \cod{\textbackslash.}
  \item \textbf{Serial} $=$ AAAALLZZNN; schimbarea lui declanșează zone transfer pe secundare.
  \item \textbf{Refresh} $\neq$ \textbf{Retry} $\neq$ \textbf{Expire}. TTL final $=$ minim(TTL, MINIMUM TTL).
\end{itemize}
\end{capcana}

\end{document}
